\chapter{Tracing an integer rasterizer}
\MARSChapterHeading{Tracing an integer rasterizer}
\label{revision-lab:integer-algorithms}

Read the main chapter's cell-selection rule, then import the student package
\texttt{007-integer-raster}. Continue its downloaded lesson in Learn. Work in the
imported \texttt{integer-line.asm}; the lesson opens the separate
\texttt{independent/dashed-line.asm} when you reach the independent exercise.
The worked source is also available at
\path{revision/labs/files/integer-line.asm} for ordinary Run and Debug outside
the guided lesson.

The inline questions ask for a prediction before execution supplies the answer.
Their source pointers identify the relevant instruction; select an answer in
Learn rather than opening Find Action. A wrong answer gives feedback and allows
another attempt. Keep a short written trace as well: a correct choice cannot
show how you obtained the next cursor, error or address.

Your evidence should connect the selected cells, the integer state and the actual
stores. A finished line can look plausible even when a tie is inconsistent or an
unrelated framebuffer word has been cleared. Test those behaviors separately.

\section{Predict cells and addresses}

Without consulting the worked trace, list the cells from $(1,1)$ to $(6,4)$.
Include both endpoints. Calculate the address of each selected cell using
$frame+4(8y+x)$. Explain why the count is six even though the horizontal difference
is five. This framebuffer has eight columns; do not reuse the paired sprite
chapter's sixteen-column stride.

Next predict $(0,0)$ to $(2,1)$. At the middle column, the continuous height is
exactly halfway between two rows. State which cell the chapter's convention
selects. Reverse the supplied endpoints and predict again. Endpoint ordering
must resolve the same geometric question in the same direction.

Choose a steep falling line, such as $(1,6)$ to $(3,0)$. State its positive
horizontal difference, negative stored vertical difference and vertical direction.
The display's downward-increasing $y$ does not change the meaning of a negative
integer step. It changes which way that step appears on the screen.

\section{Separate the snapshot from the accumulator}

Debug the default program with MIPS64 Release 6 and the Default layout. Break at
\texttt{line\_pixel}. At its first visit, identify the current point in
\texttt{\$t0}/\texttt{\$t1}, the target in \texttt{\$t2}/\texttt{\$t3}, the two
differences in \texttt{\$t4}/\texttt{\$t5}, the vertical direction in
\texttt{\$t6} and the error in \texttt{\$t7}. These values describe the cell about
to be painted, not a cell already written by this iteration.

Continue to \texttt{line\_decide}. The store and count increment have now executed,
but the doubled-error instruction has not. Step over it, record the snapshot in
\texttt{\$t8}, and predict both branches before advancing either coordinate.
After an update, verify the invariant using the new cursor and the new error.

The second visit to \texttt{line\_decide} is especially useful. At $(2,2)$, record
$e=4$ and then $e_2=8$. Step through the horizontal update: the error becomes one.
Explain why the pending vertical comparison must still read eight. Record the
next cell and its address before continuing. Compare your observations with the
main chapter's trace only after completing your prediction.

At \texttt{line\_return}, inspect the count in \texttt{\$v0} and the return link.
After returning to \texttt{line\_complete}, step over the count store. The label-addressed store is a pseudoinstruction and may require more than one
instruction step. Check \texttt{line\_count} before and after the exit-selector load. A later
\texttt{\$v0=10} does not mean that the function returned ten; it means the caller
has reused the register for another purpose.

\section{Follow the guided Debug boundaries}

Start a fresh Debug session when the lesson requests it. Use one line breakpoint
at a time, removing the previous one before resuming. A breakpoint left at
\texttt{line\_pixel} would stop on every visit, obscuring which observation the
lesson currently needs. The eight observations are:

\begin{enumerate}
  \item First \texttt{line\_pixel}, before address calculation: cursor $(1,1)$,
        error two, count zero.
  \item First \texttt{line\_decide}, before freezing the doubled error: the same
        cursor and error, count one.
  \item Second \texttt{line\_pixel}: cursor $(2,2)$, error four, count one.
  \item Second \texttt{line\_decide}: cursor $(2,2)$, error four, count two.
  \item \texttt{skip\_x} on that transition: cursor $(3,2)$, updated error one,
        saved doubled error eight, count two.
  \item Third \texttt{line\_pixel}: cursor $(3,2)$, error one, count two.
  \item \texttt{line\_return}: endpoint $(6,4)$ already painted, count six.
  \item \texttt{line\_complete}: the caller has regained control; its count store
        is still pending.
\end{enumerate}

Record whether each store is pending or complete. At the fifth observation,
explain why the vertical comparison uses eight even though the error register
now contains one. At the last observation, inspect \texttt{line\_count} before
stepping through the caller's store and exit-selector load. Finish or stop Debug
before the ordinary Run requested next. An empty console is expected; use the
bitmap and memory to inspect the raster and retained count.

The lesson saves predictions and successful checks. After editing a checked
source, continue the lesson to repeat its execution observations; an earlier
pass does not verify the edited program. The independent check uses the current
editor text, including an unsaved edit. A passing report refers to that revision.

\section{Make a fault explain itself}

In a copy, recompute \texttt{\$t8} from the current error immediately before the
vertical comparison at \texttt{skip\_x}. Predict the first cell that will differ
from the original. Then pause on the second iteration and find the instruction
that changes the decision's input. The intended error update is correct; using
that new error for an old-state decision is the fault. Repair the copy and rerun
the same counterexample.

For a different fault, make a taken horizontal update jump directly to
\texttt{line\_next}. This turns the two independent decisions into alternatives.
Use the first default transition to explain why the line now loses a required
vertical step. A diagonal movement needs both updates, not a new combined pixel
address formula.

Finally, move endpoint termination before the store. A one-point input is a short
counterexample: it still requires one painted cell and a count of one. Check the
actual store boundary rather than accepting a normal process exit as evidence
that the point was drawn.

\section{Distinguish geometry from preservation}

Test horizontal, vertical, diagonal, steep, falling and point cases. Reverse each
pair and compare all framebuffer words. A line in a different slope category
should use the same function, not a manually edited branch direction.

Prefill the frame with recognizable nonzero words before drawing. In Debug, stop
before \texttt{balc draw\_line}, edit a few words outside your predicted line, then
continue. They must retain their complete word values. Include a value with a
nonzero high byte: the display's visible RGB channels alone cannot demonstrate
preservation of every bit. The public checks automate a wider version of this
experiment with a distinct value at every cell.

Also check the endpoint record, color and guards. The function's ordering is a
change to its working registers; input memory should retain the caller's original
endpoint order. Passing the final-image checks does not justify writing beyond
the frame and later restoring a guard. The worked-program qualification therefore
observes relevant stores as well as final values.

\section{Add a dash pattern independently}

Implement \texttt{draw\_dashed} in the separate \texttt{dashed-line.asm} starter,
replacing its failing \texttt{break}. It has the same three inputs as
\texttt{draw\_line}: endpoint-record address in \texttt{\$a0}, frame base in
\texttt{\$a1} and color in \texttt{\$a2}. All four coordinates lie from zero through
seven; the frame has 64 words. Apply the same canonical endpoint and halfway rules.

Visit that raster in canonical order. Paint visit zero, preserve visit one, paint
visit two, and continue alternating. Preserve every skipped cell's original word,
every cell outside the line, the endpoint/color inputs, all guards, saved registers,
stack pointer and caller-owned stack bytes. Return the number actually painted in
\texttt{\$v0}, produce no console output and return to the supplied continuation.
The supplied caller calls twice through different continuations, storing each
returned count before continuing. Direct function checks stop at those
continuations before the caller stores anything. Preserve its \texttt{line\_count}
word during each direct call. You may use a different correct raster
algorithm; no particular temporary registers, private labels or stack layout are
required.

Before coding, distinguish three quantities: a coordinate, the number of visited
cells and the number of stores performed. Explain what each does on a skipped
cell. Using $x$ parity for the pattern fails on a vertical line. Incrementing the
result count on every visit returns the solid-line count instead of the number
painted. Neither mistake requires changing the geometric recurrence to repair it.

Use \textbf{Check this program} in Learn for feedback on the independent source.
The unfinished starter should fail; replace its \texttt{break} body and retry.
Equivalent solutions may use a different raster algorithm, private labels or
temporary registers. The public contract concerns the selected words, returned
count and preserved state, rather than a copy of the worked recurrence.

For the same public cases outside the IDE, use JDK 25 and the adjacent script:
\begin{verbatim}
java -jar /path/to/cli.jar grade \
  revision/lessons/007-integer-raster/independent/dashed-line.kts \
  /path/to/your/dashed-line.asm
\end{verbatim}
The sixteen cases vary the color, endpoint order and existing frame, including a
black line on nonzero pixels. Each case makes two calls in one machine. The
function must use its passed color, even when the global color word differs. They check the returned count, complete final words, protected
input stores and preserved function state. A correct result is accepted without
matching the worked source. Keep your predicted default dashed cells and a vertical
counterexample beside the report; a passing report alone does not explain why the
pattern follows visits rather than coordinates.


The protected endpoint record, color word and four guards must receive no
execution stores, including same-value stores or values later restored. Other
memory checks compare bytes at return; they are not a history of every access.
The checks do not prescribe the order of framebuffer stores. The caller's stack
bytes and count word must retain their values at each direct return.

Include $(0,0)$ to $(4,1)$ in your own tests. Its halfway choice occurs at visit
two, which the dashed routine paints. A wrong halfway choice at a skipped visit
could leave the final image unchanged. Compare this with $(0,0)$ to $(2,1)$,
whose halfway choice is at skipped visit one. Choose a counterexample where the
incorrect decision has an observable consequence.

Passing the sixteen public cases establishes 32 returned calls for those inputs;
it is not an exhaustive proof over all 4096 endpoint tuples, colors or previous
frames. Keep your predicted default dashed cells, vertical counterexample and
painted halfway case beside the report. Explain a failed word or count before
making another edit.
